\chapter{Find the Expensive Problem}

Myron Golden asks the aspiring entrepreneur to move attention away from a
private shortage and toward another person's problem. The shift is small in
grammar and large in consequence. \emph{I need money} describes the founder's
condition. It does not tell a customer why money should change hands.

Golden's own story includes poverty, later wealth, the loss of much of that
wealth, and a rebuilding shaped by faith. He distinguishes activity that moves
an objective from activity that merely occupies the person pursuing it. In
commercial work, the distinction becomes exact. Designing a logo, registering
a company, studying competitors, and refining a pitch can all feel like
progress. None establishes that a costly problem exists or that the proposed
work reduces it.

The first job of a business is not to express the founder. It is to make a
useful change for someone else.

That sentence sounds obvious until we ask what counts as useful, who recognizes
the change, who pays, and why the payment reaches this particular solver. A
problem can cause great human suffering and still have no functioning market.
A trivial pleasure can support a profitable market because buyers have money,
choice, and a reliable way to transact. A large market can contain no reachable
customer. A brilliant technology can arrive before anyone is willing to change
their behavior.

So we need a stricter definition of an expensive problem.

\section{Pain, cost, and payment are different}

A Houston chiropractor contrasts businesses that satisfy wants with businesses
that solve painful problems. The distinction is useful, but it should not be
turned into contempt for desire. Entertainment, beauty, hospitality, comfort,
and play create real value. A want can be frequent and important. A painful
need can be neglected because the person experiencing it cannot pay.

The chiropractor's own industry reveals why. He describes a small provider
facing concentrated insurers with greater bargaining power. The provider may
improve a patient's life while the payer determines reimbursement, rules, and
access. The social value of the work, the patient's benefit, the price paid,
and the income captured by the practitioner are four different quantities.

Tony Robbins makes a related distinction in a different context. People are
equal in dignity, he argues, while marketplaces reward contributions unequally.
That should be heard as a description of exchange, not a ranking of human
worth. Markets reward what buyers notice and can purchase under existing
institutions. Bargaining power, ownership, scarcity, law, inherited access,
timing, and external costs all affect the result. A caregiver may create
enormous value that no customer directly buys. A monopolist may capture value
partly by limiting alternatives. Revenue is evidence of exchange, not proof of
virtue.

For commercial discovery, separate four roles:

\begin{description}[leftmargin=2.4em,style=nextline]
\item[The affected person] lives with the inconvenience, loss, fear, delay, or
  unrealized possibility.
\item[The user] interacts with the proposed solution and bears its learning or
  switching cost.
\item[The decision-maker] can approve adoption, access, or a contract.
\item[The payer] supplies the budget and decides whether the exchange is worth
  its alternatives.
\end{description}

One person may occupy all four roles. In healthcare, enterprise software,
education, housing, and family purchases, they are often divided. A founder who
speaks only to the user can misunderstand the budget. A founder who speaks only
to the payer can design something the user resists. The problem is not
commercially understood until the whole decision is visible.

\section{Make the burden measurable}

An expensive problem does not need to involve a large invoice. It can be
expensive through repetition, delay, risk, attention, embarrassment, lost
revenue, unreliable quality, or the coordination required to avoid failure.
The burden can be written as a simple accounting identity:

\begin{equation}
C_{\mathrm{problem}} = C_{\mathrm{direct}} + C_{\mathrm{time}}
  + C_{\mathrm{delay}} + \mathbb{E}[C_{\mathrm{risk}}].
\end{equation}

The terms are estimates, not permission to invent precision. Direct cost is
money already spent on the problem or its workaround. Time cost is the value of
hours consumed. Delay cost is what becomes unavailable while the problem
persists. Expected risk cost combines the probability of failure with its
consequence. Some burdens---pain, dignity, anxiety, ecological damage---should
not be forced into dollars merely to make a pitch persuasive. They can remain
explicit constraints.

Frequency changes the economics. A five-minute defect encountered by ten
thousand employees each week can be more expensive than a dramatic failure
that occurs once a decade. Urgency changes the buying process. A restaurant
with a failed refrigerator needs a response now; a team mildly dissatisfied
with its reporting software can postpone a switch for years. Severity changes
the standard of proof. A tool used for medical, legal, financial, or safety
decisions must earn trust that a novelty app does not require.

The current workaround is the best starting evidence. People may use a
spreadsheet, extra staff, repeated phone calls, insurance, inventory, a
consultant, manual checking, a cheaper imperfect product, or simple endurance.
The workaround reveals what the buyer already understands, what the problem
costs in practice, and which inconvenience is apparently preferable to
switching.

Do not ask only, \emph{Would you use this?} Ask what happened the last time the
problem occurred, what the person did, how long it took, what failed, who was
involved, and what was paid. Memory of actual behavior is usually better
evidence than politeness about an imagined future.

\section{A market is a set of reachable decisions}

Neil Patel explains a simple attraction of large markets: a small captured
share of a vast pool of spending can support a substantial company. The
arithmetic is true and routinely abused.

If $n$ potential buyers spend an average of $s$ per year in a defined category,
then a rough total addressable market is

\begin{equation}
\mathrm{TAM} = n s.
\end{equation}

TAM is a ceiling under assumptions. It does not say that the buyers share one
problem, that their spending is available to a new entrant, or that one percent
is easy. The most impressive market slide often multiplies everyone who could
conceivably care by the highest imaginable price.

A useful market description narrows twice. The \emph{serviceable market}
contains buyers for whom the product, jurisdiction, language, channel, price,
and operating capacity actually fit. The \emph{obtainable market} contains the
small portion that can plausibly be reached, persuaded, served, and retained
during the next stage.

This turns market size from spectacle into a sequence of decisions:

\begin{equation}
\begin{aligned}
\text{population}&\supseteq \text{affected buyers}\\
&\supseteq \text{serviceable buyers}\\
&\supseteq \text{reachable buyers}\\
&\supseteq \text{earned customers}.
\end{aligned}
\end{equation}

A large market is valuable because it leaves room for narrow focus, not because
it excuses weak focus. The first reachable niche should be specific enough that
the same problem, language, buying process, and proof recur. A recycling
operator in Houston argues that buyers remember the leading option rather than
the second. Category leadership can help discovery, but inventing an empty
category does not create demand. The honest aim is to become unusually useful
for a narrow group, then widen only when adjacent evidence appears.

Place also changes value. A transport operator in Beverly Hills notes that the
same service can command different prices in different locations. Costs,
income, congestion, alternatives, regulation, and urgency vary geographically.
Price variation is not an invitation to exploit a captive buyer. It is evidence
that a problem has to be studied where the exchange will occur.

\section{Know the market better than the technology}

An Orlando hospitality-software executive gives the chapter's sharpest rule:
know the intended market better than the technology. The claim is deliberately
one-sided because it corrects the builder's natural bias toward the part that
can be developed without a customer in the room. It does not excuse unsafe or
weak engineering.

Market knowledge includes the user's day, the decision-maker's incentives, the
payer's budget cycle, the existing contract, the cost of changing systems, the
risk of a bad choice, the language used internally, and the event that finally
makes delay unacceptable. It includes knowing who benefits from the current
arrangement and may resist a better one.

Technology is one possible mechanism. A checklist, service, contract change,
new schedule, training program, tool, marketplace, physical object, or trusted
introduction may solve the same problem more directly. Novelty is valuable only
when it changes the outcome enough to justify adoption.

The software executive also connects research to selling: believe the case,
communicate directly, and leave a conversation with the next follow-up already
scheduled. Belief here should mean evidence strong enough to state the value
without disguise. It should not mean protecting a preferred idea from
contradictory information. A buyer who changes the founder's mind has provided
value even when no sale occurs.

This is why Golden rejects an arbitrary ninety-day demand for a million dollars
and prefers objectives tied to controllable inputs. A revenue deadline can
encourage motion and also encourage pressure, weak customers, or false claims.
A discovery objective is more honest: complete a defined number of qualified
conversations, observe the workflow, test the costly assumption, and ask for a
real commitment when the evidence is sufficient.

\section{Urgency does not suspend ethics}

Painful problems attract money because delay is costly. They also create the
conditions for exploitation. A frightened patient, desperate borrower,
isolated elder, overloaded manager, or founder near payroll may consent to
terms they would reject with time and alternatives.

An ethical problem brief therefore includes a limit on capture. What may the
solution charge, claim, collect, automate, disclose, or make difficult to leave?
Which person bears risk if the solution fails? Who lacks a voice in the
transaction? What must remain available to someone who cannot pay?

The limit is not decorative. It affects the business model. A provider may need
clear cancellation, human review, data minimization, transparent financing,
service continuity, accessibility, or referral to a better-qualified option.
The most lucrative version of a problem can be the version in which the buyer
is least able to evaluate the offer. That is precisely where restraint matters.

Nor should the founder confuse a large human need with a right to earn from it.
Permission, competence, trust, and a viable delivery system still have to be
built. Solving a severe problem badly can create more damage than leaving a
minor problem alone.

\section{Write the problem brief}

Before naming a product, write one page about the problem. Use observed facts
where they exist and label assumptions that still need testing.

\begin{enumerate}[itemsep=0.25em]
\item \textbf{Affected person.} Who experiences the problem, in what setting,
  and at which moment?
\item \textbf{Current behavior.} What happened last time, and which workaround
  did the person actually use?
\item \textbf{Burden.} What direct cost, time, delay, risk, or lost capability
  recurs, and how often?
\item \textbf{Decision.} Who uses, approves, pays, influences, and can block a
  change?
\item \textbf{Urgency.} Which event creates action, and what permits continued
  postponement?
\item \textbf{Reachable market.} Which narrow group shares the problem and can
  be reached through a credible channel now?
\item \textbf{Proof.} What observable outcome would show that the burden was
  reduced enough to justify switching and payment?
\item \textbf{Ethical limit.} Which claim, price, risk transfer, data use, or
  dependency is unacceptable even if it increases revenue?
\end{enumerate}

\enlargethispage{4\baselineskip}
The brief is not a pitch. It is a record of what is known and what would change
the builder's mind. If the affected person, workaround, burden, payer, or route
to adoption remains vague, product design is premature.

A plausible problem is still only a hypothesis. The people living with it must
be allowed to correct the description, reveal an unobserved constraint, and
show that the elegant solution solves the wrong part. The next chapter begins
with that permission to be changed.
